\setproblem
    {D}
    {컨벤션 대탈출}

\section{\pid. \pname}

\begin{frame} % No title at first slide
    \sectiontitle{\pid}{\pname}
    \sectionmeta{
        출제자: 전민주(\idtext{mingmingmon})\\
        태그: \tagtext{\#bfs, \#grid\_graph}\\
        예상 난이도: \tiertext{Silver2}
    }
    \begin{itemize}
        \item 제출 39번, 정답 13명 (정답률 33.3\%)
        \item 처음 푼 사람: 김0건 29분
    \end{itemize}
\end{frame}
%----------------------------------------------------------------------------
\begin{frame}{\textbf{\pid}. \pname}
    \begin{itemize}
        \item 주어진 컨벤션 센터의 좌석 정보에 따르면 'O'는 이동 가능한 영역, 'X'는 이동 불가능한 영역입니다.
        \item 문제에서 원하는 것은 가장 왼쪽 아래에서부터 가장 오른쪽 위까지 이동하는 최단 거리입니다.
        \item 다양한 최단경로 알고리즘(BFS, 플로이드 워샬, 데이크스트라 등) 중 적합한 것은
        한 좌석에서 다른 좌석으로 이동하는 데 항상 1만큼 이동한 것이라고 보기 때문에 BFS입니다.\\
        \item 자료 구조 Queue를 사용하여 BFS를 구현하면 해결할 수 있습니다.
        \item 이동 가능한 방향이 상,하,좌,우 4방향이므로, 격자맵을 넘어가거나,
        이미 탐색한 좌석을 재방문하지 않도록 조심해야 합니다.

    \end{itemize}
\end{frame}